\documentclass[12pt]{article}
\usepackage[utf8]{inputenc}
\usepackage[T1]{fontenc}
\usepackage{hyperref}
\title{Banda, Togbo-Vara Language Resources for K-12 Curriculum Development}
\author{Compiled by Claude Sonnet 5 (claude-sonnet-5) for Andr\'as Kornai}
\date{}
\begin{document}
\maketitle

\section{Introduction}
Togbo-Vara Banda (ISO 639-3 \texttt{tor}; Glottocode togb1241) is a Ubangian language of the Democratic Republic of the Congo and the Central African Republic. A search summary of database entries gives about 24,000 speakers. Joshua Project counts a wider people group of about 64,000 across three countries. Global Recordings Network lists two varieties, Togbo and Vara. It is distinct from Tagbu (tbm). I found no dictionary, grammar or school materials, only database entries and a mission-agency profile.

\section{Primary Educational Resources}
None found.

\section{Linguistic and Reference Materials}
\begin{itemize}
\item ASJP wordlist entry. \url{https://asjp.clld.org/languages/TOGBO_VARA_BANDA}
\end{itemize}

\section{Religious and Community Resources}
\begin{itemize}
\item Joshua Project language page. It reports that Bible translation has started and that the written status is unknown. \url{https://www.joshuaproject.net/index.php/languages/tor}
\item Joshua Project people-group page for the DRC. \url{https://joshuaproject.net/index.php/people_groups/15529/CG}
\end{itemize}

\section{Digital and Audio Resources}
\begin{itemize}
\item Global Recordings Network language page. It lists the varieties but, per the search summary, has no recordings in either one. \url{https://globalrecordings.net/en/language/tor}
\item Global Recordings Network, Togbo variety. \url{https://www.globalrecordings.net/en/language/7667}
\item Global Recordings Network, Vara variety. \url{https://globalrecordings.net/en/language/7668}
\end{itemize}

\section{Curriculum Development Context}
Nothing found supports a curriculum in this language. There is no confirmed orthography, primer or reference grammar. Any work would begin with documentation and orthography development, ideally with SIL or a local Bible-translation team. I could not verify who runs the Bible translation project.

\end{document}